%======================================================================
%  notation.tex — semantic macros.
%
%  RULE OF THUMB: never type \mathbb{E} or \mathcal{F}_t by hand again.
%  Use the macros. If you later decide that, say, quadratic variation
%  should be [X] rather than <X>, you change ONE line here and the whole
%  thesis follows. That is the entire point of this file.
%
%  The choices below resolve the collisions flagged in the review:
%    * W  is Brownian motion, always. (B is never used for it.)
%    * mu is a drift, always. Abstract measures are nu, lambda.
%    * rho is a correlation, always. Densities are p (transition), f or psi.
%    * B(t,T) is the zero-coupon bond. The letter P is left to transition
%      probabilities, correlation matrices and put prices.
%    * [0,T] is the time interval, always. (See the footnote in §1.1.)
%======================================================================

% ---------- Probability space ----------
\newcommand{\PP}{\mathbb{P}}                 % physical / real-world measure
\newcommand{\QQ}{\mathbb{Q}}                 % risk-neutral (pricing) measure
\newcommand{\QT}{\QQ^{T}}                    % T-forward measure
\newcommand{\E}{\mathbb{E}}                  % expectation
\newcommand{\Ep}[1]{\E^{#1}}                 % expectation under a named measure
\newcommand{\Var}{\mathbb{V}}                % variance
\newcommand{\Cov}{\operatorname{Cov}}
\newcommand{\Corr}{\operatorname{Corr}}
\newcommand{\Prob}{\PP}
\newcommand{\as}{\text{a.s.}}

% ---------- Sigma-algebras and filtrations ----------
\newcommand{\Fs}{\mathcal{F}}                % generic sigma-algebra
\newcommand{\Ft}{\mathcal{F}_t}              % filtration at time t
\newcommand{\Filt}{(\mathcal{F}_t)_{t\in[0,T]}}
\newcommand{\FiltW}{(\mathcal{F}^{W}_t)_{t\in[0,T]}}
\newcommand{\Gs}{\mathcal{G}}
\newcommand{\Hs}{\mathcal{H}}
\newcommand{\Borel}{\mathcal{B}}             % Borel sigma-algebra
\newcommand{\Om}{\Omega}
\newcommand{\PSpace}{(\Om,\Fs,\PP)}
\newcommand{\FPSpace}{(\Om,\Fs,\Filt,\PP)}

% ---------- Number sets ----------
\newcommand{\R}{\mathbb{R}}
\newcommand{\N}{\mathbb{N}}
\newcommand{\Z}{\mathbb{Z}}
\newcommand{\Cx}{\mathbb{C}}
\newcommand{\Tset}{\mathbb{T}}               % abstract index set (used only in §1.1)

% ---------- Processes ----------
\newcommand{\Wt}{W}                          % Brownian motion — ALWAYS W
\newcommand{\Wvec}{\mathbf{W}}               % multidimensional Brownian motion
\newcommand{\BH}{W^{H}}                      % fractional Brownian motion, Hurst H
\newcommand{\Nt}{N}                          % Poisson process
\newcommand{\qv}[1]{\langle #1 \rangle}      % predictable quadratic variation
\newcommand{\qvv}[2]{\langle #1, #2 \rangle} % predictable quadratic covariation
\newcommand{\rqv}[1]{[#1]}                   % realised (pathwise) quadratic variation
\newcommand{\rqvv}[2]{[#1,#2]}
\newcommand{\pvar}[1]{\|#1\|_{p\text{-var}}}
\newcommand{\Holder}[1]{\mathcal{C}^{#1}}

% ---------- Stochastic calculus ----------
% The differential. Chapter 1 was written with a bare italic d and Chapter 2 with
% \dd; rather than rewrite one of them, \dd is defined to BE the italic d, so the
% two chapters typeset identically. Change the line below to \mathrm{d} if the
% upright convention is ever wanted -- it will then apply to both chapters at once
% only after Chapter 1 has been converted to use the macro.
\newcommand{\dd}{d}                          % differential (italic; see note above)
\newcommand{\Lad}[1]{L^{2}_{\mathrm{ad}}\!\left(#1\right)}
\newcommand{\Lloc}[1]{\mathcal{L}\!\left(#1\right)}
\newcommand{\Gen}{\mathcal{A}}               % infinitesimal generator
\newcommand{\Genadj}{\mathcal{A}^{*}}        % its adjoint
\newcommand{\Dom}{\mathcal{D}}
\newcommand{\Expo}[1]{\mathcal{E}_{#1}}      % stochastic (Doleans-Dade) exponential
\newcommand{\Ind}{\mathds{1}}                % indicator function
\newcommand{\tr}{\operatorname{Tr}}
\newcommand{\transp}{^{\!\top}}

% ---------- Distributions ----------
\newcommand{\Normal}{\mathcal{N}}
\newcommand{\LogNormal}{\mathcal{LN}}
\newcommand{\charf}{\varphi}                 % characteristic function
\newcommand{\convd}{\xrightarrow{\;d\;}}
\newcommand{\convp}{\xrightarrow{\;\PP\;}}
\newcommand{\convLp}[1]{\xrightarrow{\;L^{#1}\;}}

% ---------- Finance ----------
\newcommand{\Spot}{S}                        % spot price
\newcommand{\Fwd}{F}                         % forward price
\newcommand{\Num}{S^{0}}                     % numeraire
\newcommand{\Bond}[2]{B(#1,#2)}              % zero-coupon bond price
\newcommand{\Disc}[2]{D(#1,#2)}              % deterministic discount factor
% The terminal payoff is G, and NOT H: the letter H is reserved for the Hurst
% parameter (\Hurst below), which is the organising symbol of Part II and appears
% already in Subsection "Stochastic Volatility Models", i.e. in the same chapter as
% the payoff. G is free throughout the financial chapters; it doubles as a generic
% event G in \mathcal{G} in the conditional-expectation material of Chapter 1, but
% the two never occur in the same chapter, let alone the same argument.
\newcommand{\Payoff}{G}                      % terminal payoff (random variable)
\newcommand{\PayoffFn}{\Phi}                 % payoff function
\newcommand{\Strike}{K}
\newcommand{\Mat}{T}                         % maturity / time horizon (deliberately shared)
\newcommand{\rf}{r}                          % risk-free short rate
\newcommand{\divy}{q}                        % continuous dividend / repo yield
\newcommand{\ivol}{\sigma_{\mathrm{imp}}}    % implied volatility
\newcommand{\lvol}{\sigma_{\mathrm{loc}}}    % local volatility
% Realised volatility: a statistic of the path under P, as opposed to \ivol, which is
% read off a price and therefore lives under Q. The three subscripts imp/loc/real are
% deliberately parallel so that the reader can tell at a glance which object is meant.
\newcommand{\rvol}{\sigma_{\mathrm{real}}}   % realised volatility
\newcommand{\tvar}{w}                        % total implied variance
\newcommand{\lmon}{k}                        % log-moneyness
% (The spot-to-forward relation is now written affinely, F = a_t S + b_t, in
%  Subsection "Dividends, and why the forward is the right variable"; no macro
%  is needed and \Theta stays free.)
\newcommand{\Snell}{U}                       % Snell envelope
\newcommand{\Stop}[1]{\mathcal{T}_{#1}}      % set of stopping times

% ---------- Rough paths (Ch. 5) ----------
%  Script capitals are unavailable (mathrsfs is not loaded, deliberately: the
%  package list is guarded for portability). Upright bold C therefore denotes
%  the SPACE of rough paths, plain \mathcal{C} the space of Holder paths, and
%  blackboard X the second level of a rough path. The three are never confused
%  because the first two carry a superscript and the third does not.
\newcommand{\RP}[1]{\mathbf{C}^{#1}}            % alpha-Holder rough paths
\newcommand{\RPg}[1]{\mathbf{C}_{g}^{#1}}       % geometric rough paths
\newcommand{\RPwg}[1]{\mathbf{C}_{wg}^{#1}}     % weakly geometric rough paths
\newcommand{\Ctrl}[1]{\mathcal{D}^{2\alpha}_{#1}} % paths controlled by #1
\newcommand{\rpX}{\mathbf{X}}                   % a rough path (X, second level)
\newcommand{\rpY}{\mathbf{Y}}
\newcommand{\lev}{\mathbb{X}}                   % its second level
\newcommand{\levW}{\mathbb{W}}                  % second level of Brownian motion
\newcommand{\Tens}[1]{T^{(#1)}}                 % truncated tensor algebra
\newcommand{\Tone}[1]{T_{1}^{(#1)}}             % its subgroup {x_0 = 1}
\newcommand{\Sig}{\mathrm{S}}                   % signature transform
\newcommand{\Sym}{\operatorname{Sym}}
\newcommand{\Anti}{\operatorname{Anti}}
\newcommand{\Hurst}{H}                          % Hurst parameter — always H
\newcommand{\rpdist}{\varrho_{\alpha}}          % rough path metric
\newcommand{\increment}[1]{\delta #1}           % the increment operator
\newcommand{\Gub}{Y'}                           % Gubinelli derivative (mnemonic)

%  Triple bar, for the homogeneous rough path norm. Defined by hand rather than
%  taken from mathtools (\DeclarePairedDelimiter) so that the package list stays
%  as it is; the spacing below is the standard one.
\newcommand{\vertiii}[1]{{\left\vert\kern-0.25ex\left\vert\kern-0.25ex\left\vert #1
    \right\vert\kern-0.25ex\right\vert\kern-0.25ex\right\vert}}

%  Shuffle product. Neither amssymb nor amsmath provides the usual symbol, so we
%  build it from \sqcup, which is what most of the literature uses in any case.
\newcommand{\shuffle}{\mathbin{\sqcup\mkern-6.5mu\sqcup}}

\makeatletter
% ---------- Editorial ----------
% Yellow markers for open items. Remove all of these before the defence.
%
% IMPORTANT: these deliberately use \colorbox and NOT soul's \hl.
% \hl re-tokenises its argument and therefore breaks on \ref, \cite, \eqref
% and most other commands ("Package soul: Reconstruction failed").
% \colorbox has no such restriction, so cross-references may be used freely
% inside the notes below. Plain \hl{...} is still fine for bare text.
%
% \figurehere{...} : full-width highlighted block, wraps over several lines.
% \todonote{...}   : same, labelled TODO. Both are block-level: \colorbox cannot
%                    break across lines, so an inline version overflows the
%                    margin whenever it lands near the end of a line.
\newcommand{\thesis@marker}[2]{%
  \par\medskip\noindent
  \colorbox{Gold}{%
    \parbox{\dimexpr\linewidth-2\fboxsep\relax}{%
      \textbf{[#1]}\ #2%
    }%
  }%
  \par\medskip
}
\newcommand{\figurehere}[1]{\thesis@marker{FIGURE HERE}{#1}}
\newcommand{\todonote}[1]{\thesis@marker{TODO}{#1}}
\makeatother

% ---------- Applications (Ch. 4): the PINN implied-volatility surface ----------
%
%  Chapter 4 is derivative-dense: the wing algebra carries three k-derivatives
%  and a T-derivative in almost every display, and the \partial_\lmon form used
%  in Chapter 2 becomes unreadable there. The partials of the surface are
%  therefore written as COMMA SUBSCRIPTS on \ivol, and \ivd is the single line
%  that fixes that convention. Reverting to \partial notation is one edit here.
%
%  The letters chosen below were forced by the collisions the notation table
%  already legislates. In particular:
%    * the core network is \Core = \mathcal{M} and NOT F, which is the forward;
%    * the analytic wing is \Wing = \mathcal{W} and NOT W, which is Brownian
%      motion always. The wing is calligraphic and always carries a side
%      subscript L or R; that is what tells the two apart, exactly as the two
%      time arguments identify the bond B(t,T).
%    * the normalised skew is \nskew = \mathfrak{s} and NOT \varphi, which is
%      the characteristic function and is needed in this very chapter for the
%      rough Heston calibration. Its relation to the wing slope s is
%      \nskew_0 = |k^*| s / a, i.e. it IS the dimensionless s; the shared
%      letter in two fonts is deliberate.
\newcommand{\ivd}[1]{\sigma_{\mathrm{imp},#1}}  % \ivd{k}, \ivd{kk}, \ivd{x}, \ivd{T}
\newcommand{\Core}{\mathcal{M}}                 % the core network (parameters \Theta)
\newcommand{\Wing}{\mathcal{W}}                 % the analytic wing
\newcommand{\Winf}{\mathcal{W}_{\infty}}        % its asymptote, a + b
\newcommand{\nskew}{\mathfrak{s}}               % normalised skew, k \ivd{k} / \ivol
\newcommand{\Sseven}{\mathcal{S}_{7}}           % the septic smoothstep
\newcommand{\Caldrift}{\Upsilon}                % the wing calendar profile
\newcommand{\elas}[1]{\lambda_{#1}}             % log-elasticity d log(#1) / d log T
